\newcommand*{\pathtwo}{chapters/chapter2}

\section{Organic Structures}


%%%%%%%%%%%%%%%%%%%%%%%%%%%%%%%%%%%%%%%%%%%%%%%%%%%%%%%%%
\subsection{Hydrocarbon frameworks and functional groups}

Organic chemistry is the study of molecules containing carbon. These molecule also contain hydrogen and sometimes oxygen, nitrogen and other elements.

\marginnote{Amino acids contain two functional groups: an amino (\ce{NH2} or \ce{NH}) group and a carboxylic acid (\ce{COOH}) group}[0cm] 
The chemistry of organic molecules depend less on the number and arrangement of hydrogen and carbon atoms than on other types of atoms. The parts of an organic molecule containing one or more other atoms are called \definition{functional groups}.
\begin{itemize}
    \item \key{The functional groups determine the way the molecule works both chemically and biologically.}
    \item \key{The hyrdocarbon framework is made up of chains and rings of carbon atoms, and it acts as a support for the functional groups.}
\end{itemize}


\subsection{Drawing molecules}

Guidelines for drawing molecules:
\begin{itemize}
    \item \key{Draw chains of atoms as zig-zags.}
    \item \key{Miss out the \ce{H}s attached to carbon atoms along with the \ce{C-H} bonds.}
    \item \key{Miss out the capital \ce{C}s representing carbon atoms.}
\end{itemize}
These are guidelines, they are not rules. If a carbon is written as \ce{C}, add all \ce{H} atoms must be written too.

It's possible to draw the same molecule in different ways on different occasions, in order to emphasize different points.


%%%%%%%%%%%%%%%%%%%%%%%%%%%%%%%%%%%
\subsection{Hyrdocarbon frameworks}
 
Carbon is unusual since it forms strong, stable bonds with the majority of elements in the periodic table (including itself).

\subsubsection*{Carbon chains}

The simplest class of hydrocarbon frameworks contains just chains of atoms (see Table \ref{tab:Cchanis}).

%\input{\pathtwo/tables/carbon_chains.tex}

\subsubsection*{Organic elements}

\marginnote{Amino acid:\\ \includegraphics[width=\marginparwidth]{\pathtwo/molecules/amino_acid}}[0cm]
\ce{R} in a structure can mean anything, is a wild card. The reactivity of organic molecules is so dependent on their functional groups that the rest of the molecule can be irrelevant and we can just choose to call it \ce{R}.

 Abbreviations for the names of carbon chains (Table \ref{tab:Cchanis}) looks like an element symbol: these symbols are sometimes called \definition{organic elements}. These symbols (and names) can be used only for terminal chains of atoms.

\subsubsection*{Carbon rings}

\marginnote{Phenyl group:\\ \includegraphics[width=\marginparwidth]{\pathtwo/molecules/phenyl_group}}[0cm] 
Rings of atoms are also common in organic structures. When a benzene ring is attached to a molecule by only one of its carbon atoms it is called a \definition{phenyl group} (with organic element symbol \ce{Ph}). Any compound containing a benzene ring or a related ring system is known as \definition{aromatic}. An \emph{aryl group} (with organic chemistry symbol \ce{Ar}) represents any substituted phenyl group, i.e. a phenyl group with any number of the hydrogen atoms replaced by other groups. Like \ce{R}, the wild card alkyl group, \ce{Ar} is a wild card aryl group.

Rings of carbon atoms are given names starting with \definition{cyclo}, following by the name for the carbon chain containing the same number of carbon atoms.

\subsubsection*{Branches}

Hydrocarbon frameworks are often branched and some short branched carbon chains are given names and organic element symbols. The most common is the isopropyl group (\ce{i-Pr}).
\begin{center}
    \includegraphics[width=\marginparwidth]{\pathtwo/molecules/isopropyl}
\end{center}

\key{Isomers are molecules with the same kinds and numbers of atoms joined up in different ways.} Isomers need not have the same functional groups.

The isobutyl group (\ce{i-Bu}) is a \ce{CH2} group joined together with an \ce{i-Pr} group. The \textit{sec}-butyl group (\ce{s-Bu}) has a methyl and an ethyl group joined to the same carbon atom. The \textit{tert}-butyl group (\ce{t-Bu}) has three methyl groups joined to the same carbon atom.

\input{\pathtwo/figures/butyl_branches.tex}

\key{The prefixes \textit{sec} and \textit{tert} are short for secondary and tertiary, terms referring to the carbon atom that attaches these groups to the rest of the molecule. A primary carbon atom is attached to only one other \ce{C} atom, a secondary to two other \ce{C} atoms, and so on.} This means there are five types of carbon atom.


%%%%%%%%%%%%%%%%%%%%%%%%%%%%%%
\subsection{Functional groups}

The understanding of functional groups is the key to the understanding of organic chemistry.

\subsubsection*{Alkanes contain no functional group}

The alkanes are the simplest lass of organic molecules because they contain no functional groups. They are extremely unreactive.

\subsubsection*{Alkenes (or olefins) contain \ce{C=C} double bonds}

\ce{C=C} double bonds impact the reactivity of an organic molecule as any other functional groups, therefore they are classified as functional groups.

\subsubsection*{Alkynes contain \ce{C#C} triple bonds}

\ce{C#C} triple bonds have a special type of reactivity associated with them, so it's useful to call a \ce{C#C} triple bond a functional group. Alkynes are linear so they are drawn with four carbon atoms in a straight line.

\subsubsection*{Alcohols (\ce{R-OH}) contain a hydroxyl (\ce{OH}) group}

Molecules containing hydroxyl groups are often soluble in water, and living organisms often attach sugar (rich in hydroxyl groups) to otherwise insoluble compounds to keep them in solution in the cell.

\subsubsection*{Ethers (\ce{R^{1}-O-R^2}) contain an alkoxy group (\ce{-OR})}

\marginnote{THF:\\ \includegraphics[width=\marginparwidth]{\pathtwo/molecules/THF}}[0cm] 
The name ethers refers to any compound that has two alkyl groups linked together through an oxygen. Tetrahydrofuran (THF) is a commonly used solvent and is a cyclic ether.

\subsubsection*{Amines (\ce{R-NH2}) contains the amino (\ce{NH2}) group}

The amino group is present in amino acids. Amines often have a powerful fishy smell. Many neurologically active compounds are also amines (like amphetamine).

\subsubsection*{Nitro compoinds (\ce{R-NO2}) contain the nitro group (\ce{NO2})}

The correct representation of the nitro group is the following:

\begin{center}
    \includegraphics[width=\marginparwidth]{\pathtwo/molecules/nitro_group}
\end{center}

Several nitro groups in one molecule can make it quite unstable and even explosive.

\subsubsection*{Alkyl halides (floruides \ce{R-F}, chlorides \ce{R-Cl}, bromides \ce{R-Br}, or iodides \ce{R-I}) contain the fluoro, chloro, bromo, or iodo groups}

These four functional groups have similar properties, although alkyl iodides are the most reactive and alkyl fluorides the least. These compounds are also known as haloalkanes (fluoroalkanes, chloroalkanes, bromoalkanes, or iodoalkanes).

\subsubsection*{Aldehydes (\ce{R-CHO}) and ketones (\ce{R^{1}-CO-R^2}) contain the carbonyl group \ce{C=O}}

\marginnote{\ce{-CHO}:\\ \includegraphics[width=\marginparwidth]{\pathtwo/molecules/aldehydes}}[0cm] 
Aldehydes can be formed by oxidizing alcohols.

\subsubsection*{Carboxylic acids (\ce{R-CO2H}) contain the carboxylic group \ce{CO2H}}

Compounds containing the carboxylic acid (\ce{CO2H}) group can react with bases, losing a proton to form carboxylate salts.

\subsubsection*{Esters (\ce{R^{1}-CO2R^2}) contain a carboxyl group with an extra alkyl group (\ce{CO2R})}

Fats are esters since they contain three esters groups. The terms ``saturated fats'' and ``unsaturated fats'' refer to whether the \ce{R} groups are saturated (no \ce{C=C} double bonds) or unsaturated (contain \ce{C=C} double bonds).

\subsubsection*{Amides (\ce{R-CONH2}, \ce{R^{1}-CONHR^2} or \ce{R^{1}-CONR^2R^3})}

Proteins are amides. They are formed when the carboxylic acid group of an amino acid condenses with the amine group of another amino acid to form an amide linkage (peptide bond).

\subsubsection*{Nitriles or cyanides (\ce{R-CN}) contain the cyano group \ce{-C=N}}

The organic nitrile group has quite different properties from those associated with lethal inorganic cyanide.

\subsubsection*{Acyl chlorides (acid chlorides, \ce{R-COCl})}

Acyl chlorides are reactive compounds used to make esters and amides. They are derivative of carboxylic acids with the \ce{-OH} replaced by \ce{-Cl}. They are too reactive to be found in nature.

\subsubsection*{Acetals}

Acetals are compounds with two single-bonded oxygen atoms to the same carbon atom.


%%%%%%%%%%%%%%%%%%%%%%%%%%%%%%%%%%%%%%%%%%%%%%%%%%%%%%%%%%%%%%%%%%%%%%%%%%%%%%%%%%%%%%%%%
\subsection{Carbon atoms carrying functional groups can be classified by oxidation level}

\marginnote{Do not confuse the oxidation level with the oxydation state. The oxidation state is determined by the number of bonds to carbon, including those to \ce{C} and \ce{H}.}[0cm] 
We say that carbon atoms carrying functional groups that can be interconverted without the need for reducing (or oxidizing) agents have the same oxidation level. \key{The oxidation level is determined by the number of heteroatoms bonded to a carbon}. \key{A heteroatom is any atom in an organic molecule other than \ce{C} or \ce{H}}.

Aldehyde and ketones contain a carbon atom with two bonds to heteroatoms; they are at the aldehyde oxidation level.

Alcohols, ethers, and alkyl halides have a carbon atom with only one single bond to a heteroatom; they are at the alcohol oxidation level.

Single alkanes, which have no bonds to heteroatoms, are at the alkane oxidation level.

Compounds that have a carbon atom with four bonds to heteroatoms is related to \ce{CO2} and best described as the carbon dioxide oxidation level.


%%%%%%%%%%%%%%%%%%%%%%%%%%%%%
\subsection{Naming Compounds}

The International Union of Pure and Applied Chemistry (IUPAC) has developed \emph{systematic nomenclature}, a set of rules that allows any compound to be given a unique name that can be deduced directly from its chemical structure. Conversely, a chemical structure can be deduced from its systematic name.

\subsubsection*{Systematic Nomenclature}

Systematic names can be divided into three parts: one describes the hydrocarbon framework, one describes the functional groups, and one indicates where the functional groups are attached to the hydrocarbon framework.

Alkyls (methyl, ethyl, propyl, ...) are the names for some simple fragments of hydrocarbon framework. Adding a hydrogen atom to these alkyl fragments and changing -yl to -ane makes the alkanes and their names.

The name of a functional group can be added to the name of a hydrocarbon framework either as a suffix or as a prefix.

\subsubsection*{Numbers are used to locate functional groups}

A number can be included in the name to indicate which carbon atom the functional group is attached to. The carbon atoms are counted from one end. When either of two numbers could be used (depending on which end you count from) and we always chose the lower of the two.

We can use use di-, tri-, and tetra- if there is more than one of the same functional group.

With cyclic compounds we can use numbers to show the distance between the two group: start from the carbon atom carrying one of the functional groups, then count round.

When hydrocarbon frameworks are branched we can name these by treating the branch as though it were a functional group.

\subsubsection*{\emph{Ortho}, \emph{meta}, and \emph{para}}

\marginnote{\emph{Ortho}, \emph{meta}, and \emph{para} are often abbreviated with \emph{o}, \emph{m}, and \emph{p}.}[0cm] 
With substituted benzene rings, an alternative way of identifying the positions of the substituents is to use the terms \emph{ortho}, \emph{meta}, and \emph{para}. \emph{Ortho} compounds are 1,2-disubstituted, \emph{meta} compounds are 1,3-disubstituted, and \emph{para} compounds are 1,4-disubstituted.

\input{\pathtwo/figures/simple_compounds.tex}

\input{\pathtwo/figures/simple_groups.tex}
